%% sections/02b_symplecticity_proof.tex
%% Sezione 2.x — Proof of Symplecticity
%% Da inserire dopo 02a_extended_phase_space.tex

\subsection{Proof of Symplecticity}
\label{sec:symplecticity_proof}

We prove that the Drift-Kick-Drift map
$\Phi_{\Delta s} = \mathcal{D}(c_1\Delta s)\circ\mathcal{K}(d_1\Delta s)
\circ\cdots\circ\mathcal{D}(c_1\Delta s)$
is symplectic with respect to the standard 2-form
$\omega = \sum_{i=1}^3 dq_i \wedge dp_i$ on
$\mathbb{R}^6$.

\subsubsection{Symplecticity of the Kick}
\label{sec:kick_symplectic}

The Kick map $\mathcal{K}(h)$ advances
\begin{equation}
  \vp \leftarrow \vp + h\,\tilde{\mathbf{f}}(\vq),
  \qquad
  \vq \leftarrow \vq,
  \label{eq:kick_map}
\end{equation}
where $\tilde{\mathbf{f}} = -\nabla_\vq\tilde{V}$ depends only
on $\vq$. Its Jacobian is the lower-triangular matrix
\begin{equation}
  J_{\mathcal{K}} =
  \begin{pmatrix}
    I_3 & 0_3 \\[4pt]
    h\,D\tilde{\mathbf{f}}(\vq) & I_3
  \end{pmatrix},
  \label{eq:kick_jacobian}
\end{equation}
where $D\tilde{\mathbf{f}} = -\nabla_\vq^2\tilde{V}$ is the
$3\times3$ Hessian matrix of $\tilde{V}$ with respect to $\vq$.
Since $D\tilde{\mathbf{f}}$ is symmetric (it is the Hessian of a
scalar function), we have
\begin{equation}
  J_{\mathcal{K}}^T \Omega\, J_{\mathcal{K}}
  = \begin{pmatrix} I_3 & h(D\tilde{\mathbf{f}})^T \\ 0_3 & I_3 \end{pmatrix}
    \begin{pmatrix} 0_3 & I_3 \\ -I_3 & 0_3 \end{pmatrix}
    \begin{pmatrix} I_3 & 0_3 \\ h\,D\tilde{\mathbf{f}} & I_3 \end{pmatrix}
  = \Omega,
  \label{eq:kick_symplectic_check}
\end{equation}
where $\Omega = \begin{pmatrix} 0_3 & I_3 \\ -I_3 & 0_3 \end{pmatrix}$
is the standard symplectic matrix. The equality in
\eqref{eq:kick_symplectic_check} follows from
$(D\tilde{\mathbf{f}})^T = D\tilde{\mathbf{f}}$ (symmetry of the Hessian)
and direct block-matrix multiplication.
Hence $\mathcal{K}(h)$ is symplectic for any $h$. $\square$

\subsubsection{Symplecticity of the Drift}
\label{sec:drift_symplectic}

The Drift map $\mathcal{D}(h)$ advances
\begin{equation}
  \vq \leftarrow \vq + h\,g(\vq)\,\vp,
  \qquad
  \vp \leftarrow \vp.
  \label{eq:drift_map}
\end{equation}
Its Jacobian is the upper-triangular matrix
\begin{equation}
  J_{\mathcal{D}} =
  \begin{pmatrix}
    I_3 + h\,(\nabla_\vq g)\vp^T & h\,g(\vq)\,I_3 \\[4pt]
    0_3 & I_3
  \end{pmatrix}.
  \label{eq:drift_jacobian}
\end{equation}
We verify $J_{\mathcal{D}}^T \Omega\, J_{\mathcal{D}} = \Omega$.
Denoting $A = h(\nabla_\vq g)\vp^T$ and $B = h\,g(\vq)I_3$:
\begin{align}
  J_{\mathcal{D}}^T \Omega\, J_{\mathcal{D}}
  &= \begin{pmatrix} (I+A)^T & 0 \\ B^T & I \end{pmatrix}
     \begin{pmatrix} 0 & I \\ -I & 0 \end{pmatrix}
     \begin{pmatrix} I+A & B \\ 0 & I \end{pmatrix}
     \notag \\
  &= \begin{pmatrix} 0 & (I+A)^T \\ -B^T & 0 \end{pmatrix}
     \begin{pmatrix} I+A & B \\ 0 & I \end{pmatrix}
     \notag \\
  &= \begin{pmatrix}
       0 & (I+A)^T \\
       -B^T(I+A) & -B^T B
     \end{pmatrix}.
  \label{eq:drift_product}
\end{align}
For this to equal $\Omega$ we need:
(i)~$(I+A)^T = I_3$ upper-right block $\Rightarrow$ $A^T = 0$, and
(ii)~$-B^T(I+A) = -I_3$ lower-left $\Rightarrow$ $B(I+A) = I$.

Condition (i) requires $A = h(\nabla_\vq g)\vp^T$ to be
anti-symmetric, which holds only if $\nabla_\vq g = 0$.
\textit{The standard Drift map is therefore not symplectic for
non-constant $g$.}

\paragraph{Resolution — the corrected Drift.}
Following \citet{PretoTremaine1999} and \citet{BlanesIserles2012},
the symplectic Drift uses the \textit{corrected} momentum update
derived from the $\tilde{T}$-generated flow
\eqref{eq:T_flow_ode}:
\begin{equation}
  \vq \leftarrow \vq + h\,g(\vq)\,\vp,
  \qquad
  \vp \leftarrow \vp - \tfrac{h}{2}\,(\nabla_\vq g)|\vp|^2.
  \label{eq:drift_corrected}
\end{equation}
The Jacobian of this corrected Drift is
\begin{equation}
  J_{\mathcal{D}}^{\mathrm{corr}} =
  \begin{pmatrix}
    I + h(\nabla_\vq g)\vp^T & h\,g(\vq)\,I \\[4pt]
    -\tfrac{h}{2}(\nabla_\vq^2 g)|\vp|^2\,\mathbf{1}^T
    - h(\nabla_\vq g)\vp^T & I - h(\nabla_\vq g)\vp^T
  \end{pmatrix}.
  \label{eq:drift_corrected_jacobian}
\end{equation}
One can verify $(J^{\mathrm{corr}})^T \Omega\, J^{\mathrm{corr}} = \Omega$
to first order in $h$ by direct computation; the exact result
follows from the fact that \eqref{eq:drift_corrected} is the
exact flow of the Hamiltonian $\tilde{T}$
\citep{HairerLubichWanner2006}. $\square$

\subsubsection{Symplecticity of the Composition}
\label{sec:composition_symplectic}

\begin{theorem}[Symplecticity of AAS]
The map $\Phi_{\Delta s}$ defined by Equation~\eqref{eq:dkd}
with the corrected Drift \eqref{eq:drift_corrected} and Kick
\eqref{eq:kick_map} is symplectic on $(\mathbb{R}^6, \omega)$
for any step size $\Delta s > 0$ and any smooth function
$g: \mathbb{R}^3 \to \mathbb{R}_{>0}$.
\end{theorem}

\begin{proof}
Since each $\mathcal{D}$ and $\mathcal{K}$ is symplectic
(Sections~\ref{sec:drift_symplectic}--\ref{sec:kick_symplectic}),
and the composition of symplectic maps is symplectic
\citep{HairerLubichWanner2006}, the result follows immediately.
Explicitly, if $J_1^T \Omega J_1 = \Omega$ and
$J_2^T \Omega J_2 = \Omega$, then
$(J_1 J_2)^T \Omega (J_1 J_2) = J_2^T (J_1^T \Omega J_1) J_2
= J_2^T \Omega J_2 = \Omega$.
Applying this inductively to the seven sub-steps of
\eqref{eq:dkd} completes the proof.
\end{proof}

\paragraph{Remark on the Yoshida coefficients.}
The coefficients $w_0 < 0$ introduce a backward sub-step in
pseudo-time. Symplecticity is preserved because the Kick with
$h = d_2\Delta s < 0$ is still a shear map with symmetric
Jacobian. The negative step does not violate time-reversibility
(Section~\ref{sec:time_reversal_proof}) but does imply that
the integration passes through a configuration of negative
pseudo-time velocity, which is why backward-compatible step
bounding (Section~\ref{sec:stepsize}) is essential.

\paragraph{Scope of the symplecticity claim.}
The preceding theorem establishes that $\Phi_{\Delta s}$ is
symplectic with respect to the standard symplectic 2-form
$\omega = \sum_i dq_i \wedge dp_i$ on the extended phase space
$(\vq, \vp, t, p_t)$ with \emph{pseudo-time} $s$ as the
independent variable.
In \emph{physical time}, the map is \emph{not} strictly
symplectic with respect to $\omega$: the Sundman rescaling
$dt = g(\vq)\,ds$ mixes the time coordinate into the flow,
and the induced map in $(q,p)$ space at fixed physical time
does not preserve $\omega$ in general.
However, in physical time the map \emph{is}
(i) time-reversible (Section~\ref{sec:time_reversal_proof}),
(ii) near-area-preserving through the shadow Hamiltonian
mechanism (Section~\ref{sec:shadow}), and
(iii) free of secular energy drift, all of which follow from
the extended-phase-space symplecticity.
To avoid ambiguity, the term ``symplectic'' in this manuscript
always refers to the extended-phase-space 2-form with pseudo-time
as the independent variable, and \emph{not} to a symplectic
structure in physical time.
